\section{Question 1}

\begin{enumerate}
    \item Bootstrap resampling was used on the training data from the \emph{insurance.xlsx} dataset provided in assignment 1 question 1. This was done to estimate the uncertainty of the
          model intercept. This was done on both the linear regression model and the logistic regression model used in assignment 1 question 1. This process was done by firstly using
          the same preprocessing pipeline and model setup that was used in that question. Thereafter, a loop of a 1000 iterations was setup in which the model was fit on training data
          that was resampled with replacement using the \emph{resample} method from scikit-learn (bootstrapping). To ensure variability, the random\_state parameter was set to the current
          iteration number in the \emph{resample} method. After the model was trained, the intercept value was obtained from the model and stored. To develop a 95\% confidence
          interval of the intercept of each model, the 2.5\% and 97.5\% percentiles were obtained of the generated bootstrap distribution of the intercept. This can be done since the
          distribution mimics how the actual statistic would behave if you had pulled completely new samples in a real world scenario. The results can be seen below in 
          Figure \ref{fig:bootstrap_linear} and Figure \ref{fig:bootstrap_logistic}:

          \begin{figure}[h!]
            \centering
            \includegraphics[width=0.55\textwidth]{images/bootstrap_linear.png}
            \caption{Bootstrap Distribution of the Linear Regression Model Intercept}
            \label{fig:bootstrap_linear}
          \end{figure}

          \begin{figure}[h!]
            \centering
            \includegraphics[width=0.55\textwidth]{images/bootstrap_logistic.png}
            \caption{Bootstrap Distribution of the Logistic Regression Model Intercept}
            \label{fig:bootstrap_logistic}
          \end{figure}

    \item Thereafter, the intervals from question 1.1 was used to test the hypothesis that the intercept is equal to 0. At a 5\% significance level, one can use the 95\% bootstrap
          confidence interval since this is a two-sided test. The following hypothesis test was setup:

          Null Hypothesis:
          \[H_0: \beta_0 = 0\]

          Alternative Hypothesis:
          \[H_1: \beta_0 \neq 0\]

          The rule:
          \[\text{Reject H$_0$ if 0} \notin CI_{95\%}\]
          \[\text{Fail to reject H$_0$ if 0} \in CI_{95\%}\]

          Based on this setup, one would reject the null hypothesis if 0 does not lie within the 95\% confidence interval.

          From Figure \ref{fig:bootstrap_linear}, the 95\% bootstrap confidence interval of the Linear Regression model's intercept is [-0.2419, 0.3503]. Based on this, 0 lies
          within the lower and upper bounds of this interval. Therefore, we fail to reject H$_0$ at a 5\% significance level. Thus, there is insufficient evidence to conclude
          that the Linear Regression model's intercept is significantly different from 0.

          From Figure \ref{fig:bootstrap_logistic}, the 95\% bootstrap confidence interval of the Logistic Regression model's intercept is [-0.8130, -0.7100]. Based on this, 
          0 does not lie within the lower and upper bounds of this interval. Therefore, we reject H$_0$ at a 5\% significance level. Thus, there is sufficient evidence to conclude
          that the Logistic Regression model's intercept is significantly different from 0.

\end{enumerate}
